\section{AI 项目的风险-收益空间}

讲者借用投资管理中的经典概念——\textbf{Risk}（风险）和 \textbf{Reward}（收益）——构建了一个分析 AI 项目的二维框架。

\subsection{风险-收益矩阵}

在这个框架中，横轴表示 Reward（收益/实用性），纵轴表示 Risk（风险/潜在危害）。讲者与学生互动，将不同 AI 应用标注在这个空间中：

\begin{table}[H]
\centering
\caption{AI 应用在风险-收益空间中的定位}
\begin{tabular}{@{}lccl@{}}
\toprule
\textbf{AI 应用} & \textbf{Risk} & \textbf{Reward} & \textbf{说明} \\
\midrule
Code Completion（单行） & 低 & 低 & 填充函数签名，错了影响不大 \\
Text Revision（文本修订） & 低--中 & 中 & 给回改进版文本，风险略高 \\
Text Summary（文本摘要） & 中 & 中--高 & 可能遗漏或误解关键信息 \\
LLM（大语言模型） & 中--高 & 高 & 能力惊人但存在多种风险 \\
AI Judge（自动法官） & 很高 & 不确定 & 涉及人的自由与生命 \\
AGI（通用人工智能） & 极高 & 极高 & 高收益伴随不可控风险 \\
\bottomrule
\end{tabular}
\end{table}

\begin{warningbox}{AGI 的双重极端性}
讲者将 AGI 定位在风险-收益空间的右上角——既有极高的潜在收益（通用智能可解决人类面临的几乎所有问题），也有极高的风险（不可预测的行为、潜在的失控）。这种``高风险-高收益''的组合使得 AGI 成为最具争议的 AI 目标。如同电影《终结者》中的场景，超级智能一旦失控，后果将无法挽回。
\end{warningbox}

\subsection{风险阈值：在哪里画线？}

讲者提出了一个关键问题：\textbf{我们应该在风险坐标轴的哪个位置画一条``不可逾越''的红线？}

\begin{itemize}
    \item \textbf{保守策略}：红线画得很低，只允许低风险应用（如代码补全、文本修订）
    \item \textbf{激进策略}：红线画得很高，容忍大部分风险，只禁止极端场景
    \item \textbf{现实困境}：不同利益相关方对风险阈值的判断截然不同
\end{itemize}

有学生提到了一个重要的第三维度：\textbf{经济实力}。大公司有充足的法律团队应对风险后果，而初创公司可能无力承受法律纠纷的成本。这种不对称性使得``画线''问题更加复杂。

\begin{knowledgebox}{法律只关心风险，不关心收益}
讲者指出一个重要区别：法律体系\textbf{不考虑收益维度}。你不能在法庭上辩护说``我的 AI 很有用所以应该被允许违规''——法律只关心你的系统是否造成了伤害或违反了规定。这意味着在法律框架下，风险-收益分析被简化为纯粹的风险评估。
\end{knowledgebox}

\subsection{本章小结}

风险-收益框架为思考 AI 项目提供了一个直观的分析工具。不同的 AI 应用在这个空间中占据不同位置，而``在哪里画线''本质上是一个涉及技术、伦理、法律和经济的多维决策问题。法律体系倾向于只关注风险维度，而开发者需要在收益诱惑和风险约束之间找到平衡点。

